\documentclass[10pt]{extarticle}
\usepackage[a4paper,margin=1.5cm]{geometry}
\usepackage[T1]{fontenc}
\usepackage{amsmath,amssymb}
\usepackage{graphicx}
\usepackage[export]{adjustbox}
\usepackage{multicol}
\usepackage{enumitem}
\usepackage{booktabs,array,tabularx}
\setlength{\columnsep}{0.7cm}
\setlength{\parindent}{0pt}
\setlength{\parskip}{3pt}
\pagestyle{empty}
\makeatletter
% figure registry (figdefs.tex is generated by the extract step)
\newcommand{\qfigdef}[4]{\@namedef{qf@file@#1}{#2}\@namedef{qf@w@#1}{#3}\@namedef{qf@wide@#1}{#4}}
\newcommand{\qf@check}[1]{\@ifundefined{qf@file@#1}{\PackageError{score2}{Unknown figure id #1}{}}{}}
\newcommand{\qf@img}[4][]{\includegraphics[width=#2\dimexpr\@nameuse{qf@w@#4}pt\relax,max width=#3,max height=0.45\textheight,#1]{\@nameuse{qf@file@#4}}}
% Full-width questions: finalize wraps every question that contains a wide figure in \qfullbegin...\qfullend,
% so the WHOLE question (stem, figures, options) is set across the page in reading order, then the two columns
% resume. The question's head - its number, stem and first figure (\qfullsplit marks the end of the head) -
% is kept on one page, so a number is never left at a page foot with its figure overleaf; the rest may break
% across pages normally, so no big gaps are left. Without a \qfullsplit, the whole question is kept together
% if it is under 45% of a page. Never write these yourself.
\newif\ifqf@full
\newif\ifqf@boxed
\newsavebox{\qf@box}
\newcommand{\qf@open}{\global\qf@fulltrue\par\global\qf@boxedtrue
  \setbox\qf@box\vbox\bgroup\hsize\textwidth\linewidth\textwidth}
\newcommand{\qf@emit}[1]{\par\egroup\global\qf@boxedfalse
  \ifdim\dimexpr\ht\qf@box+\dp\qf@box\relax<#1\textheight\noindent\usebox{\qf@box}\par\else\unvbox\qf@box\fi}
\newcommand{\qfullsplit}{\ifqf@boxed\qf@emit{0.85}\fi}
\newcommand{\qf@close}{\ifqf@boxed\qf@emit{0.45}\fi\global\qf@fullfalse}
\newcommand{\qfullbegin}{\end{multicols}\qf@open}
\newcommand{\qfullend}{\qf@close\begin{multicols}{2}}
\newcommand{\qfullstart}{\qf@open}
\newcommand{\qfulldone}{\qf@close}
% \qfig[scale]{id}: block figure, centred, directly where it is written (size set by the tool).
% Fallback if a wide figure is not inside a full-width question: break out of the columns in place - no floats.
\newcommand{\qfig}[2][1]{\qf@check{#2}%
  \ifqf@full\qf@blk{#1}{#2}\else
  \ifnum\@nameuse{qf@wide@#2}=1 \expandafter\@firstoftwo\else\expandafter\@secondoftwo\fi
  {\end{multicols}\begin{center}\qf@img{#1}{\textwidth}{#2}\end{center}\begin{multicols}{2}}%
  {\qf@blk{#1}{#2}}\fi}
\newcommand{\qf@blk}[2]{\par\smallskip{\centering\qf@img{#1}{\linewidth}{#2}\par}\smallskip}
% \qfigi[scale]{id}: inline figure (inside option lists / tables / minipages); never breaks columns.
\newcommand{\qfigi}[2][1]{\qf@check{#2}\qf@img[valign=c]{#1}{\linewidth}{#2}}
\makeatother
% structure macros - lint reads these to check completeness and figure placement
\newcommand{\papersec}[1]{\par\bigskip{\centering\Large\bfseries PAPER-#1\par}\medskip}
\newcommand{\partsec}[2]{\par\medskip\noindent{\large\bfseries PART-#1: #2}\par\smallskip}
\newcommand{\sect}[1]{\par\smallskip\noindent\textbf{#1}\par}
\newcommand{\qn}[1]{\par\smallskip\noindent\textbf{#1.}~}
\newcommand{\paraq}[2]{\par\smallskip\noindent\textbf{Paragraph for Questions #1 and #2}\par}
\newenvironment{qoptions}{\begin{itemize}[leftmargin=1.8em,itemsep=1pt,topsep=2pt,label={}]}{\end{itemize}}
\newcommand{\opt}[2]{\item[(#1)] #2}
% figure registry (generated by score2_tools.py finalize)
\qfigdef{p1_phy_s0_q01_1}{SRGBT - Test 09 (19 Sep 2021)-img/p1_phy_s0_q01_1.png}{78.6}{0}
\qfigdef{p1_phy_s0_q02_1}{SRGBT - Test 09 (19 Sep 2021)-img/p1_phy_s0_q02_1.png}{110.0}{0}
\qfigdef{p1_phy_s0_q03_1}{SRGBT - Test 09 (19 Sep 2021)-img/p1_phy_s0_q03_1.png}{156.6}{0}
\qfigdef{p1_phy_s0_q04_1}{SRGBT - Test 09 (19 Sep 2021)-img/p1_phy_s0_q04_1.png}{94.8}{0}
\qfigdef{p1_phy_s0_q05_1}{SRGBT - Test 09 (19 Sep 2021)-img/p1_phy_s0_q05_1.png}{53.5}{0}
\qfigdef{p1_phy_s0_q06_1}{SRGBT - Test 09 (19 Sep 2021)-img/p1_phy_s0_q06_1.png}{87.1}{0}
\qfigdef{p1_phy_s0_q01_2}{SRGBT - Test 09 (19 Sep 2021)-img/p1_phy_s0_q01_2.png}{77.6}{0}
\qfigdef{p1_phy_s0_q08_1}{SRGBT - Test 09 (19 Sep 2021)-img/p1_phy_s0_q08_1.png}{81.0}{0}
\qfigdef{p1_phy_s0_q09_1}{SRGBT - Test 09 (19 Sep 2021)-img/p1_phy_s0_q09_1.png}{147.8}{0}
\qfigdef{p1_phy_s0_q10_1}{SRGBT - Test 09 (19 Sep 2021)-img/p1_phy_s0_q10_1.png}{213.1}{0}
\qfigdef{p1_phy_s0_q11_1}{SRGBT - Test 09 (19 Sep 2021)-img/p1_phy_s0_q11_1.png}{140.2}{0}
\qfigdef{p1_phy_s0_q12_1}{SRGBT - Test 09 (19 Sep 2021)-img/p1_phy_s0_q12_1.png}{113.1}{0}
\qfigdef{p1_phy_s0_q13_1}{SRGBT - Test 09 (19 Sep 2021)-img/p1_phy_s0_q13_1.png}{123.5}{0}
\qfigdef{p1_phy_s0_q14_1}{SRGBT - Test 09 (19 Sep 2021)-img/p1_phy_s0_q14_1.png}{121.6}{0}
\qfigdef{p1_phy_s0_q15_1}{SRGBT - Test 09 (19 Sep 2021)-img/p1_phy_s0_q15_1.png}{61.3}{0}
\qfigdef{p1_phy_s0_q16_1}{SRGBT - Test 09 (19 Sep 2021)-img/p1_phy_s0_q16_1.png}{95.0}{0}
\qfigdef{p1_phy_s0_q17_1}{SRGBT - Test 09 (19 Sep 2021)-img/p1_phy_s0_q17_1.png}{121.6}{0}
\qfigdef{p1_phy_s0_q18_1}{SRGBT - Test 09 (19 Sep 2021)-img/p1_phy_s0_q18_1.png}{74.1}{0}
\qfigdef{p1_chem_s0_q25_1}{SRGBT - Test 09 (19 Sep 2021)-img/p1_chem_s0_q25_1.png}{181.2}{0}
\qfigdef{p1_chem_s0_q28_1}{SRGBT - Test 09 (19 Sep 2021)-img/p1_chem_s0_q28_1.png}{181.3}{0}
\qfigdef{p1_chem_s0_q31_1}{SRGBT - Test 09 (19 Sep 2021)-img/p1_chem_s0_q31_1.png}{124.7}{0}
\qfigdef{p1_chem_s0_q33_1}{SRGBT - Test 09 (19 Sep 2021)-img/p1_chem_s0_q33_1.png}{202.7}{0}
\qfigdef{p1_chem_s0_q34_1}{SRGBT - Test 09 (19 Sep 2021)-img/p1_chem_s0_q34_1.png}{207.2}{0}
\raggedcolumns
\makeatletter\@beginparpenalty=10000\makeatother
\newcolumntype{L}[1]{>{\raggedright\arraybackslash}p{#1}}
\begin{document}
\begin{center}{\Large\bfseries SRGBT --- Test 09 (19 Sep 2021)}\end{center}
\begin{multicols}{2}
\raggedcolumns
\papersec{1}
\partsec{1}{PHYSICS}
\qn{1} In the figure is a string attached at one end to wall and other end to block B of mass 1kg. The total length of string is 35 cm but the vertical part is of length 10 cm. B is resting against vertical wall of A of mass 3 kg free to move on the floor. The minimum value of coefficient of static friction between A and floor to avoid slipping anywhere is [Neglect size of pulley]\par
\qfig{p1_phy_s0_q01_1}
\begin{qoptions}
\opt{A}{$\dfrac15$}
\opt{B}{$\dfrac37$}
\opt{C}{$\dfrac4{17}$}
\opt{D}{$\dfrac5{12}$}
\end{qoptions}
\qn{2} The system shown is in equilibrium, positive value of x$_1$ and x$_2$ denote extensions of the springs, while negative values denote compressions. Which of the following is not possible ?\par
\qfig{p1_phy_s0_q02_1}
\begin{qoptions}
\opt{A}{x$_1$ = 3 cm, x$_2$ = 2 cm}
\opt{B}{x$_1$ = 8 cm, x$_2$ = 3 cm}
\opt{C}{x$_1$ = 1 cm, x$_2$ = $-$8 cm}
\opt{D}{x$_1$ = $-$1 cm, x$_2$ = $-$2 cm}
\end{qoptions}
\qn{3} $\mu_\mathrm{AB}$ = 0.5; $\mu_\mathrm{CD}$ = 0.8\par
$\mu_{\mathrm B,\text{ table}}$ = $\mu_{\mathrm D,\text{ table}}$ = 0\par
\qfig{p1_phy_s0_q03_1}
The magnitude of friction between C and D is :
\begin{qoptions}
\opt{A}{80 N}
\opt{B}{40 N}
\opt{C}{50 N}
\opt{D}{60 N}
\end{qoptions}
\qn{4} A small metal ball is being pulled gradually on a fixed frictionless hemisphere as shown in the figure. Radii of the ball and that of the pulley are much smaller than the radius of the hemisphere. As the ball slides from the bottom to a position close to the top of the hemisphere, how do the magnitudes of pulling force F and contact force R between the ball and the hemisphere change?\par
\qfig{p1_phy_s0_q04_1}
\begin{qoptions}
\opt{A}{F increases and R decreases.}
\opt{B}{F decreases and R increases.}
\opt{C}{F decreases and R remains unchanged.}
\opt{D}{F remains unchanged and R decreases.}
\end{qoptions}
\qn{5} Consider the situation shown in the adjacent figure and select the correct statement(s) from the following.\par
\qfig{p1_phy_s0_q05_1}
\begin{qoptions}
\opt{A}{Acceleration of pulley B is g/11.}
\opt{B}{Acceleration of pulley C is g/11.}
\opt{C}{Tension in the string passing over pulley A is 12.46 g.}
\opt{D}{Tension in the string passing over pulley A is 10 g.}
\end{qoptions}
\qn{6} Coefficient of friction between the two blocks Q and R is 0.3 whereas the surface AB is smooth\par
\qfig{p1_phy_s0_q06_1}
\begin{qoptions}
\opt{A}{Acceleration of P is 5.86 m/s$^2$}
\opt{B}{Tension T$_1$ in the string is 17.7 N}
\opt{C}{Tension T$_2$ in the string is about 41.4 N}
\opt{D}{Acceleration of P is 7.55 m/s$^2$}
\end{qoptions}
\qn{7} Two blocks having same mass are placed on rough incline plane and the coefficient of friction between A and incline is $\mu_1$ = 1.0 and between block B and incline is $\mu_2=\dfrac34$. As the inclination of the plane `$\theta$' with respect to horizontal increases, choose the correct answer (s).\par
\qfig{p1_phy_s0_q01_2}
\begin{qoptions}
\opt{A}{There is no contact force between block A and B as $\theta\le37^\circ$}
\opt{B}{There is no contact force between block A and B as $\theta\le45^\circ$}
\opt{C}{As $\mu_1>\mu_2$ both the blocks move together}
\opt{D}{They starts moving at an angle $\theta=\tan^{-1}\dfrac78$}
\end{qoptions}
\qn{8}
\qfig{p1_phy_s0_q08_1}
$\mu_\mathrm{AB}$ = 0.5; $\mu_\mathrm{BC}$ = 0.2; $\mu_{\mathrm C,\mathrm{ground}}$ = 0.1\par
Choose the correct statement(s)
\begin{qoptions}
\opt{A}{The acceleration of A is 3 m/s$^2$}
\opt{B}{The acceleration of C is 1 m/s$^2$}
\opt{C}{The acceleration of C is 3 m/s$^2$}
\opt{D}{The force of friction between A and B is 20 N.}
\end{qoptions}
\qn{9} Consider ideal pulleys and string. The accelerations of B and C are given as 2 and 3 m/s$^2$ respectively. Choose the correct statements :\par
\qfig{p1_phy_s0_q09_1}
\begin{qoptions}
\opt{A}{The magnitude of acceleration of A is $\sqrt{181}$ m/s$^2$}
\opt{B}{The magnitude of acceleration of A is $\sqrt{141}$ m/s$^2$}
\opt{C}{The acceleration of point P w.r.t. ground of the string is 7 m/s$^2$}
\opt{D}{The acceleration of point P w.r.t. ground of the string is 5 m/s$^2$}
\end{qoptions}
\qn{10}
\qfig{p1_phy_s0_q10_1}
The acceleration of A is 3 m/s$^2(-\hat{\mathrm j})$. Choose the correct statement(s).
\begin{qoptions}
\opt{A}{The acceleration of B is 4$\hat{\mathrm i}$ m/s$^2$}
\opt{B}{The acceleration of B is $-4\hat{\mathrm i}$ m/s$^2$}
\opt{C}{The acceleration of C is $\left[\left(\dfrac{64}{25}\right)\hat{\mathrm i}+\left(\dfrac{48}{25}\right)\hat{\mathrm j}\right]$m / s$^2$}
\opt{D}{The acceleration of C is $\left[\left(\dfrac{64}7\right)\hat{\mathrm i}+\left(\dfrac{48}7\right)\hat{\mathrm j}\right]$m / s$^2$}
\end{qoptions}
\qn{11}
\qfig{p1_phy_s0_q11_1}
$\mu_\mathrm{AB}$ = 0.8; $\mu_\mathrm{BC}$ = 0.5; $\mu_{\mathrm C,\text{ table}}$ = 0\par
The maximum possible value of mass of block D so that A, B, C move together without any relative motion among themselves is found to be K times 5 kg. The value of K is \underline{\hspace{1.5cm}} .
\qn{12} A plank of mass 10 kg rests on a smooth horizontal surface. Two blocks A and B of masses m$_\mathrm A$ = 2 kg and m$_\mathrm B$ = 1 kg lies at a distance of 3 m on the plank as shown in the figure. The friction coefficient between the blocks and plank are $\mu_\mathrm A$ = 0.3 and $\mu_\mathrm B$ = 0.1. Now a force F = 15 N is applied to the plank in horizontal direction. Find the time after which block A collides with B.\par
\qfig{p1_phy_s0_q12_1}
\qn{13} Three blocks A, B and C each of mass m are arranged in pulley string system as shown in the figure. All pulleys are massless and frictionless, while all strings are light and inextensible. Coefficient of kinetic friction, $\mu$ between block A and horizontal surface is equal to 0.5 and force P acts on it. The value of P (in Newton) which block C does not move, is (m = 0.75 kg)\par
\qfig{p1_phy_s0_q13_1}
\qn{14} Three blocks A, B, C of masses m, m and 2m respectively consists a system. Neglect friction. After the system is released, at any instant `t$_0$' the speed of block `C' w.r.t ground becomes $(\sqrt3+1)$ m/s.\par
m$_\mathrm A$ = m$_\mathrm B$ = m; m$_\mathrm C$ = 2m; (neglect friction)\par
The speed of block B w.r.t C at given instant `t$_0$' is\dots\dots.(in m/s)\par
\qfig{p1_phy_s0_q14_1}
\qn{15} On the diagram, two blocks of equal mass are connected by an ideal string. The values of m = 1kg and k$_2$ = 100 N/m are given. Initially, both springs are relaxed. Then the left block is slowly pulled down a distance 0.1 m and released. Find maximum possible value of k$_1$ (in N/m) for which both the blocks will have same magnitude of acceleration just after releasing. Both spring have natural length of 2m.\par
\qfig{p1_phy_s0_q15_1}
\qn{16}
\qfig{p1_phy_s0_q16_1}
$\mu_\mathrm{AB}$ = 0.5; $\mu_\mathrm{BC}$ = 0.2; $\mu_{\mathrm C,\mathrm{ground}}$ = 0.1\par
The acceleration of A is \underline{\hspace{1.2cm}} m/s$^2$.
\qn{17} At t = 0, block of mass m = 5 kg is given an initial velocity of 12 m/s towards right and a constant force of 5N always acts towards left as shown in figure. Cofficient of friction between block and ground is $\mu$ = 0.2. Distance covered by block till t = 5 sec is \underline{\hspace{1.2cm}} m.\par
\qfig{p1_phy_s0_q17_1}
\qn{18} An object of mass m = $\sqrt8$ kg rests on an inclined plane that makes angle $\theta$ = 45$^\circ$ with the horizontal floor. What minimum force (in N) , parallel to the base of the incline must be applied to the object in order to begin to move it along the plane parallel to the floor as shown? The coefficient of static friction between the object and the plane is $\mu_\mathrm s$ = 1.25.\par
\qfig{p1_phy_s0_q18_1}
\raggedcolumns
\papersec{1}
\partsec{2}{CHEMISTRY}
\qn{19} A mixture of NH$_3$(g) and N$_2$H$_4$(g) is placed in a sealed container at 300 K. The total pressure is 0.5 atm. The container is heated to 1200 K at which both substances decompose completely according to equations :\par
2NH$_3$ $\longrightarrow$ N$_2$ + 3H$_2$ and N$_2$H$_4$ $\longrightarrow$ N$_2$ + 2H$_2$\par
After decomposition is complete, the total pressure at 1200 K is found to be 4.5 atm. Find the mole percent of each gas in the original mixture.
\begin{qoptions}
\opt{A}{25\% NH$_3$, 75\% N$_2$H$_4$}
\opt{B}{75\% NH$_3$, 25\% N$_2$H$_4$}
\opt{C}{82\% NH$_3$, 18\% N$_2$H$_4$}
\opt{D}{80\% NH$_3$, 20\% N$_2$H$_4$}
\end{qoptions}
\qn{20} The radial part of schrodinger wave equation for hydrogen atom is\par
$\psi(\mathrm r)=\dfrac1{16\sqrt4\,\mathrm a_0^{3/2}}\,(\sigma-1)\,(\sigma^2-8\sigma+12)\,\mathrm e^{-\sigma/2}$\par
Where a$_0$ = constant \& $\sigma$ = 2r/na$_0$; n = principle quantum number\par
Select the incorrect statement:
\begin{qoptions}
\opt{A}{Distance of nearest radial node from the nucleus is 2a$_0$.}
\opt{B}{Distance of farthest radial node from the nucleus is 12a$_0$.}
\opt{C}{Number of maxima in the curve $\psi$(r) vs r are 4.}
\opt{D}{$\psi$(r) is for 4p orbital.}
\end{qoptions}
\qn{21} Two vessel `A' and `B' containing same mass of H$_2$(g) and CH$_4$(g) respectively. A is maintained at 300 K and `B' is maintained at 600 K. Volume of A and B are same. Gases behave ideally, then identify incorrect statement.\par
\textbf{Given :} $\left(\text{radius, r}_{\mathrm{CH}_4}=2\mathrm r_{\mathrm H_2}\right)$
\begin{qoptions}
\opt{A}{Pressure of gas in A is more than pressure in B}
\opt{B}{Average speed in A is more than B}
\opt{C}{Mean free path of gas in A is less than mean free path in B}
\opt{D}{Pressure in both A and B are same}
\end{qoptions}
\qn{22} The angular part of the wave function depends on the quantum numbers
\begin{qoptions}
\opt{A}{n, $\ell$, m}
\opt{B}{n, $\ell$}
\opt{C}{$\ell$, m}
\opt{D}{only m}
\end{qoptions}
\qn{23} Identify the correct statements is/are :
\begin{qoptions}
\opt{A}{A real gas `X' with dominating intermolecular force of attraction will have slower rate of diffusion as compared to same gas behaving ideally under identical condition.}
\opt{B}{CH$_4$ diffuses at faster rate as compared to O$_2$ gas (gases behave ideally and gases are under identical condition).}
\opt{C}{At absolute zero, thermal motion of the molecules of gases cease.}
\opt{D}{Density of moist air is less than dry air at same temperature \& pressure (gases behave ideally \& air mainly contains oxygen \& nitrogen).}
\end{qoptions}
\qn{24} A particle of unknown charge \& mass is projected towards center of copper atom with an initial speed of x m/s from a large distance the closest distance of approach is observed to be a$_0$ meters. Using the information, identify the correct option(s).\par
[K = 9 $\times$ 10$^9$ N m$^2$/col$^2$ ; e = 1.6 $\times$ 10$^{-19}$ col]
\begin{qoptions}
\opt{A}{The specific charge of particle is given by $\dfrac{1\mathrm x^2\times\mathrm a_0}{2\times29\mathrm{Ke}}$}
\opt{B}{If initial speed is halved, closest distance of approach will become half of original.}
\opt{C}{If initial speed is halved, closest distance of approach will become four times of original.}
\opt{D}{If distance between particle is greater than a$_0$ then possible speed can be $\le$ x m/s.}
\end{qoptions}
\qn{25}
\qfig{p1_chem_s0_q25_1}
If 6 moles of ``A'' are introduced into the vessel through the value, following reaction ocurs at isothermal conditions.\par
3A$_{\mathrm{(g)}}$+B$_{\mathrm{(g)}}$ $\to$ AB$_{\mathrm{(s)}}$ + A$_{2\mathrm{(g)}}$\par
Select incorrect statement among following :
\begin{qoptions}
\opt{A}{No limiting reagent in this reaction}
\opt{B}{Product gases exert total pressure = 3 atm}
\opt{C}{Final $\Delta$h value = 76cm}
\opt{D}{If liquid (X) is used in tube insted of Hg (density of X = 2$\times$density of Hg), then final $\Delta$h = 152cm}
\end{qoptions}
\qn{26} Identify the statements which is/are \textbf{incorrect} w.r.t model of an atom.
\begin{qoptions}
\opt{A}{Probability density at the nucleus for s and p orbital depends on ``n'' \& ``n, m'' respectively as per wave mechanical model.}
\opt{B}{As per Bohr's theory ground state angular momentum is equal to $\dfrac{\mathrm h}{2\pi}$.}
\opt{C}{The distance of antinode in the radial probability distribution graph of 1s orbital for H, He$^+$, Li$^{2+}$ is in the ratio 6 : 3 : 2.}
\opt{D}{As per wave mechanical model; penetration power as well as angular momentum increases as the value of azimuthal quantum number increases.}
\end{qoptions}
\qn{27} Select the correct statement(s) :
\begin{qoptions}
\opt{A}{At critical conditions volume occupied by the gas is 12 times the voume of 1 mole gaseous molecules.}
\opt{B}{A gas can be liquified at its T$_\mathrm b$ (Boyle's temperature) by application of pressure.}
\opt{C}{At very high pressure and moderate temperature, the volume of n mole real gas is equal to $\dfrac{\mathrm{nRT}}{\mathrm P}$ + nb.}
\opt{D}{An ideal gas can never be liquified.}
\end{qoptions}
\qn{28} An electron is ejected from Li$^{2+}$ on exposing 600 nm, the stopping potential \& Wavelength of ejected electron is:\par
(Given : h = 6.626 $\times$ 10$^{-34}$ Js, m$_\mathrm e$ = 9.1 $\times$ 10$^{-31}$ kg)\par
\qfig{p1_chem_s0_q28_1}
\begin{qoptions}
\opt{A}{1.675 eV, 3.83 $\times$ 10$^{-10}$ nm}
\opt{B}{1.6625 eV, 1.204 nm}
\opt{C}{2.66 $\times$ 10$^{-19}$ J, 1.204 nm}
\opt{D}{6.4 $\times$ 10$^{-20}$ J, 3.83 $\times$ 10$^{-10}$ nm}
\end{qoptions}
\qn{29} A 2 litre rigid container is present inside a large 6 litre container. 2 litre container contain 12 mole of H$_2$ \& 10 mole of He. While 6 litre container 20 mole of He and 18 mole of CH$_4$. There is an orifice on one wall of small container which can be operated by a remote control device. Calculate total number of moles present inside small container when orifice is opened for sufficiently long time.
\qn{30} (i) \quad The Schrodinger wave equation for hydrogen atom is $\Psi_{2s}=\dfrac1{4\sqrt{2\pi}}\left(\dfrac1{\mathrm a_0}\right)^{3/2}\left(2-\dfrac{\mathrm r_0}{\mathrm a_0}\right)\mathrm e^{-\frac{\mathrm r_0}{\mathrm a_0}}$, where a$_0$ is Bohr's radius. If the radial node for 2s be at r$_0$, then the value of $\dfrac{\mathrm r_0}{\mathrm a_0}$ is `\textbf{X}'.\par
(ii) \quad The atomic masses of He and Ne are 4 and 20 amu, respectively. The value of the de--Broglie wavelength of He gas at $-73^\circ$C is `M' times that of the de-Broglie wavelength of Ne at 727$^\circ$C. The value of M is `\textbf{Y}'.\par
\textbf{Find the value of X $\times$ Y.}
\qn{31} The container I is found to contain 2 mole H$_2$, 4 mole Ne and 4 mole CH$_4$ when the connecting tube is in opened position.\par
\qfig{p1_chem_s0_q31_1}
(a) \quad What is the total number of moles of gases in container maintained at 1200 K.\par
(b) \quad Partial pressure of H$_2$ in container `I' [R = 0.08 L-atm K$^{-1}$mol$^{-1}$]\par
Give your answer $\dfrac{(\mathrm a)\times(\mathrm b)}{10}$.
\qn{32} A sample containting large number of Li$^{2+}$ ions all in a particular excited state is subjected to a photons coming from M shell to K shell transition of any unielectronic species. If the sample is further excited to a higher level \& de--excites back to the ground state then calculate energy of the photon (eV) having 8$^\mathrm{th}$ highest energy among all emitted from de--excitation.
\qn{33} Select correct statements in context to graphical information given for 2 substances `A' \& `B'. (a, b are Vander Waal's parameter)\par
\qfig{p1_chem_s0_q33_1}
(i) \quad Normal boiling point of liquid `A' is less than 300 K, if curve - x is almost rectangular hyperbolic.\par
(ii) \quad If Boyle's temperature for gas `A' is 250 K, then gas `A' has Z $>$ 1 at T = 300 K.\par
(iii) \quad If critical temperature for gas `B' is 450 K, then $\left(\dfrac{\mathrm V_1}4<\mathrm b\right)$ and its P$_\mathrm{Critical}$ $>$ P$_1$.\par
(iv) \quad Critical temperature of both gases `A' \& `B' can never be same.\par
(v) \quad Vapour pressure of liquid `B' at 400 K is P$_1$.
\qn{34} A sample containing large number of atoms of unielectronic species is subjected to EMR having wavelength = $\dfrac{1240}{10.2}$nm cuasing excitation in the sample which contains some atoms in ground state as well as some in one particular excited state; ``A''. On de--excitation to the ground level total 6 different wavelengths were observed out of which 3 wavelength were shorter than $\dfrac{1240}{10.2}$nm . From this info answer the questions that follow.\par
[Use : hc = 1240 eV-nm; I.E. of hydrogen = 13.6 eV/atom.\par
(i) The value of ``n'' for the original excited state ``A'' ?\par
(ii) The atomic number of the unielectronic species involved.\par
\textbf{Give your answer as (i) $\times$ (ii)}
\qn{35}
\qfig{p1_chem_s0_q34_1}
(i) Value `x' remains closed as piston reaches B \quad (ii) Value `x' remains closed as piston reaches C\par
(iii) Value `y' remains closed as piston reaches B \quad (iv) Value `y' remains closed as piston reaches C\par
(v) P$_\mathrm{total}$ = 25 atm as piston reaches B \quad (vi) P$_\mathrm{total}$ = 50 atm as piston reaches C.\par
If no. of correct statements = x\par
If no. of incorrect statements = y\par
then $\dfrac{\mathrm x}{4\mathrm y}$ value = ?
\qn{36} A gas of identical H-like atom has some atoms in the lowest (ground) energy level A and few atoms in a particular upper (excited) energy level B and there are no atoms in any other energy level. The atoms of the gas make transition to a higher energy level by absorbing monochromatic light of photon energy 2.7 eV. Subsequently, the atoms emit radiation of only six different photons energies. Some of the emitted photons have energy 2.7 eV. Some have more and some have less than 2.7 eV. How many statements is/are correct ?\par
(i) \quad The principal quantum number of initially excited level B = 2.\par
(ii) \quad The ionisation energy for the gas atoms is 28.8 eV.\par
(iii) \quad The maximum and the minimum energies of the emitted photons will be 13.5 eV and 0.699 eV respectively.\par
(iv) \quad In above question nucleus is stationary \& e$^-$ will revolve around the nucleus in circular path.\par
(v) \quad If all emitted photons have energy $\ge$ 2.7 eV, then in this case principal quantum number of initially excited level B = 3.
\raggedcolumns
\papersec{1}
\partsec{3}{MATHEMATICS}
\qn{37} Let $\alpha,\beta$ be the roots of $x^2+px+1=0$ and $\gamma,\delta$ be the roots of $x^2+qx+1=0$. What is the value of following expression?\par
P=$(\alpha-\gamma)(\beta-\gamma)(\alpha+\delta)(\beta+\delta)$
\begin{qoptions}
\opt{A}{$q^2-p^2$}
\opt{B}{$p^2-q^2$}
\opt{C}{$(p+q)^2$}
\opt{D}{$(p-q)^2$}
\end{qoptions}
\qn{38} What is the minimum value of $f(x)=\dfrac{(x+a)(x+b)}{(x+c)}$, where $x>-c$, $a>c$, $b>c$ :
\begin{qoptions}
\opt{A}{$\sqrt{a-c}+\sqrt{b-c}$}
\opt{B}{$\left(\sqrt{a-c}+\sqrt{b-c}\right)^2$}
\opt{C}{$2\left(\sqrt{a-c}+\sqrt{b-c}\right)$}
\opt{D}{$2\left(\sqrt{a-c}+\sqrt{b-c}\right)^2$}
\end{qoptions}
\qn{39} Let f(x) = Ax$^2$ + Bx + C (A, B, C $\in$ R) if f(x) is an integer whenever x is an integer, then 2A, A+B and C are all integers.
\begin{qoptions}
\opt{A}{statement is always true}
\opt{B}{statement is never true}
\opt{C}{statement is true in some cases not always}
\opt{D}{none of these}
\end{qoptions}
\qn{40} If $\alpha$ is a real root of the quadratic equation $ax^2+bx+c=0$ and $\beta$ is a real root of $-ax^2+bx+c=0$, then
\begin{qoptions}
\opt{A}{there is a root $\gamma$ of of the equation $\dfrac a2x^2+bx+c=0$ which lies between $\alpha$ and $\beta$}
\opt{B}{both root of $\dfrac a2x^2+bx+c=0$ lies between $\alpha$ and $\beta$}
\opt{C}{no root of $\dfrac a2x^2+bx+c=0$ lies between $\alpha$ and $\beta$}
\opt{D}{data is insufficient}
\end{qoptions}
\qn{41} If $6x^5+5x^3+x^2+4x-1=a(x-1)^5+b(x-1)^4+c(x-1)^3+d(x-1)^2+e(x-1)+f$, then
\begin{qoptions}
\opt{A}{a + 2b +3c + 4d + 5e + 6f = 910}
\opt{B}{a + 2b +3c + 4d + 5e + 6f $\ne$ 920}
\opt{C}{a $-$ b + c $-$ d + e $-$ f = 1}
\opt{D}{a $-$ b + c $-$ d + e + f = 2}
\end{qoptions}
\qn{42} The possible integers $\lambda$ such that $(x-\lambda)(x-10)+1=0$ has integer roots is :
\begin{qoptions}
\opt{A}{12}
\opt{B}{8}
\opt{C}{20}
\opt{D}{24}
\end{qoptions}
\qn{43} $\sqrt{2x^2+5x-2}-\sqrt{2x^2+5x-9}=1$, possible value of x is/are :
\begin{qoptions}
\opt{A}{2}
\opt{B}{3}
\opt{C}{$-\dfrac92$}
\opt{D}{$\dfrac92$}
\end{qoptions}
\qn{44} If real numbers x such that $\left(x-\dfrac1x\right)^{\frac12}+\left(1-\dfrac1x\right)^{\frac12}=x$ is $\dfrac{\sqrt a+1}b$, then
\begin{qoptions}
\opt{A}{a = 5}
\opt{B}{b = 2}
\opt{C}{a + b = 7}
\opt{D}{ab = 5}
\end{qoptions}
\qn{45} Let $f(x)=ax^2+bx+c$, where a, b, c $\in$ R, suppose that $|f(x)|\le1\ \forall\ x\in[0,1]$. Which of the following statemetns are true?
\begin{qoptions}
\opt{A}{$|$ a $|\le$ 8}
\opt{B}{$|$ b $|\le$ 8}
\opt{C}{$|$ c $|\le$ 1}
\opt{D}{$|$ a $|+|$ b $|+|$ c $|\le17$}
\end{qoptions}
\qn{46} If $\alpha,\beta,\gamma,\delta$ are the roots of $120x^4-134x^3-73x^2+4x+3=0$, $|\alpha|<|\beta|<|\gamma|<|\delta|$ then
\begin{qoptions}
\opt{A}{$(\alpha.\beta.\gamma)^{-1}$ is equal to 60}
\opt{B}{$(\alpha+\beta+\gamma+\delta)$ is equal to $\dfrac{134}{120}$}
\opt{C}{all four roots are rational}
\opt{D}{two roots are rational and two are irrational}
\end{qoptions}
\qn{47} The value of k for which $x^4+(1-2k)x^2+(k^2-1)=0$ has three real solution is :
\qn{48} If a$_1$, a$_2$\dots..a$_{2\mathrm n}$ be in AP. If $\mathrm a_1^{\,2}-\mathrm a_2^{\,2}+\mathrm a_3^{\,2}-\mathrm a_4^{\,2}+\dots\dots+\mathrm a_{2\mathrm n-1}^2-\mathrm a_{2\mathrm n}^2=\dfrac{\mathrm n}{\alpha\mathrm n-1}\left(\mathrm a_1^{\,\beta}-\mathrm a_{2\mathrm n}^{\,\gamma}\right)$, then $\alpha+\beta+\gamma$ is:
\qn{49} Let \quad $x=1+3a+6a^2+10a^3+\dots\dots\infty$, $|a|<1$\par
$y=1+4b+10b^2+10b^3+\dots\dots\infty$, $|b|<1$\par
If \quad $1+3(ab)+5(ab)^2+\dots\dots\infty=\dfrac{1+(1-x^\alpha)(1-y^\beta)}{\left[1-(1-x^\gamma)(1-y^\delta)\right]^2}$ then $(\alpha\,\beta\,\gamma\,\delta)^{-1}$
\qn{50} Minimum value of $10x^2+5y^2-4x-6y-12yx+18$ is
\qn{51} If $\displaystyle\sum_{r=1}^{n}\mathrm t_r=\dfrac18\mathrm n(\mathrm n+1)(\mathrm n+2)(\mathrm n+3)$, then If $\displaystyle\sum_{r=1}^{n}\dfrac1{\mathrm t_r}=\dfrac12-\dfrac{\mathrm a}{(\mathrm n+\mathrm b)(\mathrm n+\mathrm c)}$, then a + b + c is :
\qn{52} If $\displaystyle\sum_{m=1}^{\infty}\sum_{n=1}^{\infty}\dfrac{\mathrm m^2\mathrm n}{3^\mathrm m(\mathrm n3^\mathrm m+\mathrm m3^\mathrm n)}=\dfrac{\mathrm p}{\mathrm q}$, where p and q are coprime then q $-$ p is:
\qn{53} $\dfrac14+\dfrac{1.3}{4.6}+\dfrac{1.3.5}{4.6.8}+\dots\dots$ is equal to:
\qn{54} The ninth term of 3, 5, 10, 20, 39 is :
\end{multicols}
\end{document}
